This document describes how I carried out the first two months of my capstone on glucose forecasting for type 1 diabetes, from early September to late October. The project began as a replication of GLIMMER, a published study that trains forecasters to weight errors in dangerous glucose ranges more heavily. A bug in that replication erased its main gain, and my first uncertainty bands raised far too many false alarms. Those turns led to the question I study now: whether the band around a 60-minute forecast stays honest in each glucose situation, for each patient and on each day. I explain how I came to define success, how a written research protocol and a register that keeps every failed hypothesis guide my testing, and how I checked results on eight outside studies never used to design the methods. I also describe my planning against what actually happened, my use of an AI coding assistant and the checks I place on it, how I used my advisor's feedback, and the work that remains before the Full Draft.